\documentclass[10pt,a4paper]{article}
\usepackage{amsmath,amssymb,amsthm,graphicx}
% Minimal retained preamble for lane-p2 cards taken from
% sources/ahl-16-fixed-points-2019/editorial-preamble.tex
% (doi:10.5802/ahl.16, Annales Henri Lebesgue 2 (2019) 149-185, CC BY 4.0).
%
% Copied verbatim from the editorial preamble: the package list, the
% theorem declarations under the shared counter and the macro block, so that
% every macro and environment of the retained spans is available.
%
% Omitted and why: the class line (the assembler supplies the class); the
% a4paper geometry line and the fontspec block with its explicit lmroman10
% files (encodings and fonts fight the forced xelatex engine, and the
% assembler supplies the page setup).
\usepackage{amsmath,amssymb,amsthm,mathtools,graphicx,xcolor,hyperref,xurl,xspace,ifthen}
\newtheorem{theo}{Theorem}[section]
\newtheorem{lemm}[theo]{Lemma}
\newtheorem{theorem}[theo]{Theorem}
\newtheorem*{remas*}{Remarks}
\providecommand{\og}{«\nobreakspace}
\providecommand{\fg}{\nobreakspace»}
\newtheorem{lem}[theo]{Lemma}
\newtheorem{lemma}[theo]{Lemma}
\newtheorem{prop}[theo]{Proposition}
\newtheorem{coro}[theo]{Corollary}
\newtheorem{corol}[theo]{Corollary}
\newtheorem{conj}[theo]{Conjecture}
\newtheorem{Proposition}[theo]{Proposition}
\newtheorem{theoreme}[theo]{Theorem}
\theoremstyle{definition}
\newtheorem{defi}[theo]{Definition}
\newtheorem{exam}[theo]{Example}
\newtheorem{exem}[theo]{Example}
\newtheorem{rema}[theo]{Remark}
\newtheorem{nota}[theo]{Notation}
\newtheorem*{rema*}{Remark}
\newtheorem*{defi*}{Definition}
\newtheorem*{theo*}{Theorem}
\newtheorem*{lemm*}{Lemma}
\newtheorem*{prop*}{Proposition}
\newcommand{\enoncename}{}
\newtheorem{corpusenonce}[theo]{\enoncename}
\newenvironment{enonce}[2][]{\renewcommand{\enoncename}{#2}\begin{corpusenonce}}{\end{corpusenonce}}
\newenvironment{enonce*}[2][]{\par\medskip\noindent\textbf{#2.}\itshape}{\par\medskip}
\newenvironment{proofc}[1][\proofname]{\begin{proof}[#1]}{\end{proof}}
\newcommand{\Mc}{\mathcal{M}}
\newcommand{\pointir}{.\ ---\ }
\providecommand{\mathof}[1]{\mathrm{#1}}

\numberwithin{equation}{section}
\numberwithin{figure}{section}
\numberwithin{table}{section}
\def\endappendix{}
%~Mouliné par MaN_auto v.0.21.4 2019-04-30 09:02:09


\usepackage[all]{xy}

\hyphenation{auto-morphism
every
Roelcke
peri-od-ic
peri-od-icy
rep-re-sen-ta-tion
}








\newcommand{\A}{\mathcal A}
\newcommand{\N}{\mathbb{N}}
\newcommand{\Z}{\mathbb{Z}}
\newcommand{\Q}{\mathbb{Q}}
\newcommand{\K}{\mathcal{K}}
\newcommand{\CC}{\mathbb{C}}

\newcommand*{\mA}{\mathbf{A}}
\newcommand*{\mB}{\mathbf{B}}
\newcommand*{\mC}{\mathbf{C}}
\newcommand*{\mD}{\mathbf{D}}
\newcommand*{\mE}{\mathbf{E}}
\newcommand*{\mF}{\mathbf{F}}
\newcommand*{\mH}{\mathbf{H}}
\newcommand*{\mZ}{\mathbf{Z}}
%\newcommand*{\mA}{\mathbf{A}}
%\newcommand*{\mA}{\mathbf{A}}

\newcommand{\restrict}[2]{#1 \upharpoonright #2}

\newcommand{\e}{\varepsilon}

\DeclareMathOperator{\Aut}{Aut}
\DeclareMathOperator{\Stab}{Stab}
\DeclareMathOperator{\Age}{Age}

\newcommand{\actson}{\curvearrowright}

\newcommand{\RUCB}{\mathrm{RUC}_{b}}
\newcommand{\ROB}{\mathrm{Ro}_{b}}
\newcommand{\WAP}{\mathrm{WAP}}
\newcommand{\AP}{\mathrm{AP}}

\graphicspath{{./figures/}}

\newcommand*{\mk}{\mkern -1mu}
\newcommand*{\Mk}{\mkern -2mu}
\newcommand*{\mK}{\mkern 1mu}
\newcommand*{\MK}{\mkern 2mu}

\hypersetup{urlcolor=purple, linkcolor=blue, citecolor=red}

\newcommand*{\romanenumi}{\renewcommand*{\theenumi}{\roman{enumi}}}
\newcommand*{\Romanenumi}{\renewcommand*{\theenumi}{\Roman{enumi}}}
\newcommand*{\alphenumi}{\renewcommand*{\theenumi}{\alph{enumi}}}
\newcommand*{\Alphenumi}{\renewcommand*{\theenumi}{\Alph{enumi}}}
\let\oldexists\exists
\renewcommand*{\exists}{\mathrel{\oldexists}}
\let\oldforall\forall
\renewcommand*{\forall}{\mathrel{\oldforall}}


\usepackage{cleveref}
\crefname{originalproposition}{original proposition}{original propositions}
\Crefname{originalproposition}{Original proposition}{Original propositions}
\makeatletter
\newlabel{prop:RoelckeRP}{{4.5}{source}{}{sourcepointer-1}{}}
\newlabel{prop:RoelckeRP@cref}{{[originalproposition][0][]4.5}{[1][0][]source}{}{}{}}
\makeatother
\crefname{originalproposition}{original proposition}{original propositions}
\Crefname{originalproposition}{Original proposition}{Original propositions}
\makeatletter
\newlabel{prop:Roelcke coloring}{{4.4}{source}{}{sourcepointer-2}{}}
\newlabel{prop:Roelcke coloring@cref}{{[originalproposition][0][]4.4}{[1][0][]source}{}{}{}}
\makeatother
\title{The Ramsey property for Roelcke colorings is equivalent to its one-point joint embedding form}
\author{Lionel Nguyen Van Thé}
\date{doi:10.5802/ahl.16; CC BY 4.0}
\begin{document}
\maketitle
\noindent\textit{Editorial note.} The counted unit is the article's own Proposition prop:RPRoelcke at source lines 527-529 together with its proof at source lines 531-533, a straight adjacent pairing in the source file. The source is the publisher version of the article in Annales Henri Lebesgue 2 (2019), pages 149-185, DOI 10.5802/ahl.16, released under CC BY 4.0.\par\medskip
\noindent\textit{Editorial note.} The two neighbouring statements of the same section are not counted. The one at lines 519-521 carries the same equivalence with a stronger transfer condition and has its own proof at lines 523-525; the one at lines 515-517 is the membership criterion for a single coloring with its own proof at lines 523-525. This card counts the statement at lines 527-529, whose proof at lines 531-533 is its own block: the forward direction is a specialisation of the quantified form, and the converse is the ultrahomogeneity lifting, which is the substantive step and is carried out here in full.\par\medskip
\noindent\textit{Editorial note.} The licence evidence used here is the observation file kept with the source, license-observation.json, which records the licence as CC BY 4.0, points at the inherited publisher-version evidence and states that the landing page displays no licence line; the underlying evidence files are the same ones the earlier cards of this source rely on. No acceptance record is written anywhere.\par\medskip
\noindent\textit{Editorial note.} Both pointers of this card carry the number the article prints for its statements: its theorem-like environments share one counter attached to the section, so the printed number is the section together with the place of the statement inside it, read off an enumeration of every numbered theorem-like environment of the source in order, which gives 4.5 and 4.4. The anchors recorded point at each statement's own lines.\par\medskip
\noindent\textit{Editorial note.} The gate raised one manual-review note: the next proof environment begins at line 535. That block belongs to the following statement of the section and is not part of this card; the proof retained here is closed at line 533.\par\medskip
\noindent\textit{Editorial note.} Neither the retained statement nor the retained proof contains a citation, so this card needs no bibliography entry.\par\medskip
\noindent\textit{Editorial note.} The census package marked this candidate free of missing citation keys and listed the statement span as 527-529 and the proof span as 531-533; both were re-read from the source before the card was written.\par\medskip
\noindent\textit{Editorial note.} rqual sv-mismatch batch fix 2026-10-14 (Chief ruling: re-assemble unpublished only): source\_version aligned verbatim to the source dir license-observation.json sv (package evidence authority); revision bumped for re-assembly. Recipe from lane-s3.\par\medskip
\noindent\textit{Allowed standard prerequisites (editorial).}\begin{itemize}
\item Roelcke colorings and joint embedding patterns. A finite coloring of the copies of a finite substructure is a Roelcke coloring when it depends only on the joint embedding pattern of the copy with a fixed finite substructure of the structure, and the joint embedding patterns are the objects the article uses to index those values, as made precise by its definitions. The author's own definitions and the surrounding results are the ones used by the counted statement and proof.. Both the statement and the proof of the counted proposition are phrased entirely in these objects: the one-point form asks for a single joint embedding of that shape whose associated coloring is constant on the required copies.
\end{itemize}\medskip
\ifdefined\hypertarget\hypertarget{sourcepointer-1}{}\fi
\noindent\textit{Source reference.} Original proposition 4.5: sources/ahl-16-fixed-points-2019/body-01.tex, source lines 519--521. Navigation pointer invoked by the retained proof at its start, where the quantified form of the same equivalence supplies the copy in the forward direction. That proposition is another statement of the same section and is not retained in this card, so the reference is declared as a pointer carrying the number the article itself prints for that statement, that is the section and the place of the statement inside that section, counted by enumerating every numbered theorem-like environment of the source in order: it is 4.5; the cited statement is an allowed input and none of it is reproduced inside the span.\par\medskip
\ifdefined\hypertarget\hypertarget{sourcepointer-2}{}\fi
\noindent\textit{Source reference.} Original proposition 4.4: sources/ahl-16-fixed-points-2019/body-01.tex, source lines 515--517. Navigation pointer invoked by the retained proof only inside the description of the earlier quantified form, whose hypothesis speaks of colorings attached to joint embedding patterns. That statement is the membership criterion for a single Roelcke coloring, an earlier result of the section, and it is not retained in this card, so the reference is declared as a pointer carrying the number the article itself prints for that statement, counted over the numbered theorem-like environments of the source in order: it is 4.4; the cited statement is an allowed input and none of it is reproduced inside the span.\par\medskip
\begin{prop}\label{prop:RPRoelcke}
Let $\mF$ be a Fra\"{\i}ss\'e structure. Then $\mF$ has the Ramsey property for colorings in $\ROB(\Aut(\mF))$ iff for every $\mA, \mB, \mZ\in \Age(\mF)$, and for every finite coloring $\gamma$ of the joint embedding patterns of $\mA$ and $\mZ$, there exists a joint embedding $\langle b,z\rangle$ such that the coloring $a\mapsto \gamma([a,z])$ is constant on $\binom{b(\mB)}{\mA}$.
\end{prop}

\begin{proof}
Assume that the Ramsey property for colorings in $\ROB(\Aut(\mF))$ holds, and fix $\mA, \mB, \mZ\in \Age(\mF)$, $\gamma$ a finite coloring of the joint embedding patterns of $\mA$ and $\mZ$. Fix $z\in \binom{\mF}{\mZ}$. Then $b\in \binom{\mF}{\mB}$ obtained by Proposition~\ref{prop:RoelckeRP} is as required. Conversely, fix $\mA, \mB, \mZ\in \Age(\mF)$, $z\in \binom{\mF}{\mZ}$, $\gamma$ a finite coloring of the joint embedding patterns of $\mA$ and $\mZ$. Consider a joint embedding $\langle b',z'\rangle$ such that $\gamma([\cdot,z'])$ is constant on $\binom{b'(\mB)}{\mA}$. Let $i$ be the unique isomorphism such that $i\circ z'=z$. Then by ultrahomogeneity of $\mF$, we can extend $i$ to $b'(\mB)\cup z'(\mZ)$, and $b:=i\circ b'$ is as required.
\end{proof}

\end{document}
